%
% PhD thesis style
% 
% Improved by Vincent X. Gong
% with love, October 2019
%
% Updated: Jul 2021
% http://www.gong.im
% http://https://www.gong.im/phd_thesis_latex_trail_template/
%
% originally created by Rudy Negenborn
%
% based on trail thesis style for LaTeX
%
% February 2008
%
% further info: rudy@negenborn.net
% and: http://www.negenborn.net/mampc/
%

% How to get started? For a start:
% * read through this file
% * edit 0-misc/defs.tex
% * edit 0-misc/title.tex
% * latex thesis.tex
% and there you have your first version

\documentclass[12pt,hidelinks]{trailthesis}
%\documentclass[a4paper,12pt]{book}
\usepackage{etex}
\usepackage{geometry}
\geometry{left=1cm,right=1cm}

% include the draft option to include tags in the .dvi showing were
% overfull hboxes are located, e.g.,
% \documentclass[10pt,draft]{trailthesis}

% include the hidelinks option to remove the red box in the table of content

%%%%%%%%%%%%%%%%%%%%%%%%%%%%%%%%%%%%%%%%%%%%%%%%%%%%%%%%%%%%%%%%

% define the base directory in which your thesis files are located
%
%
% NOTE WINDOWS USERS: the '\' in the Windows directory name should
%                     be replaced by '/', so:
%                     C:/thesis/, and not C:\thesis\
%
% NOTE WINDOWS USERS: make sure that you thesis files are in
%                     a directory *without* spaces in the name,
%                     so: C:/thesis, and not C:/My Document/My Thesis
%
% NOTE: make sure that the directory name ends with a '/'
%
%\newcommand{\basedir}{/Users/vgong/Documents/workspaces/phd_thesis/thesis_trail_vg/}

%\usepackage[xetex]{hyperref}
%\usepackage{hyperref}
%%%%%%%%%%%%%%%%%%%%%%%%%%%%%%%%%%%%%%%%%%%%%%%%%%%%%%%%%%%%%%%

% include following package to be able to have line numbers in from of
% all text of thesis. this can be useful to include in the draft
% version of your thesis, such that the committee members can easily
% point out where they have remarks/suggestions/corrections.

\usepackage[left,modulo]{lineno}

%%%%%%%%%%%%%%%%%%%%%%%%%%%%%%%%%%%%%%%%%%%%%%%%%%%%%%%%%%%%%%%%

\newcommand\dedication[1]{%
	\thispagestyle{empty}%
	\vspace*{\fill}%
	\begin{flushright}%
		#1%
	\end{flushright}%
%	\vspace*{\fill}%
%	\cleardoublepage%
}
\newcommand\epigraph[3][2\baselineskip]{%
	\begin{flushright}%
		{\rmfamily\itshape #2}%
		\vskip 0.5\baselineskip%
		#3%
	\end{flushright}%
	\vspace{#1}%
}



% include base definitions used in this thesis
%%%%%%%%%%%%%%%%%%%%%%%%%%%%%%%%%%%%%%%%%%%%%%%%%%%%%%%%%%%%%

% define note command

\setlength{\marginparwidth}{30mm}

\newcommand{\mynote}[1]{%
\marginpar{\footnotesize{#1}}%
}%


% uncomment to remove remarks.
%\renewcommand{\mynote}[1]{} 

%%%%%%%%%%%%%%%%%%%%%%%%%%%%%%%%%%%%%%%%%%%%%%%%%%%%%%%%%%%%%

% define psfrag centering command
\newcommand{\psfragc}[2]{%
\psfrag{#1}[t][t]{\color[rgb]{0,0,0}%
\setlength{\tabcolsep}{0pt}%
\begin{tabular}{c}#2\end{tabular}}%
}


%%%%%%%%%%%%%%%%%%%%%%%%%%%%%%%%%%%%%%%%%%%%%%%%%%%%%%%%%%%%%
% define the bib style

\usepackage[authoryear]{natbib}
%\usepackage[]{natbib}

%%%%%%%%%%%%%%%%%%%%%%%%%%%%%%%%%%%%%%%%%%%%%%%%%%%%%%%%%%%%%

% adjust maximum number of columns that can be used in a matrix environment

% add matrix cols
%\setcounter{MaxMatrixCols}{30}

%%%%%%%%%%%%%%%%%%%%%%%%%%%%%%%%%%%%%%%%%%%%%%%%%%%%%%%%%%%%%
% all about spaces

\raggedbottom % Prevents spreading out vertical space
\parskip=6pt % Fixed space between paragraphs
\baselineskip=12pt % Fixed baseline skip


%%%%%%%%%%%%%%%%%%%%%%%%%%%%%%%%%%%%%%%%%%%%%%%%%%%%%%%%%%%%%
% config the dedication section

\newcommand\dedication[1]{%
	\thispagestyle{empty}%
	\vspace*{\fill}%
	\begin{flushright}%
		#1%
	\end{flushright}%
%	\vspace*{\fill}%
%	\cleardoublepage%
}
\newcommand\epigraph[3][2\baselineskip]{%
	\begin{flushright}%
		{\rmfamily\itshape #2}%
		\vskip 0.5\baselineskip%
		#3%
	\end{flushright}%
	\vspace{#1}%
}


%%%%%%%%%%%%%%%%%%%%%%%%%%%%%%%%%%%%%%%%%%%%%%%%%%%%%%%%%%%%%%

% define note command

\setlength{\marginparwidth}{30mm}

\newcommand{\mynote}[1]{%
\marginpar{\footnotesize{#1}}%
}%


% uncomment to remove remarks.
%\renewcommand{\mynote}[1]{} 

%%%%%%%%%%%%%%%%%%%%%%%%%%%%%%%%%%%%%%%%%%%%%%%%%%%%%%%%%%%%%

% define psfrag centering command
\newcommand{\psfragc}[2]{%
\psfrag{#1}[t][t]{\color[rgb]{0,0,0}%
\setlength{\tabcolsep}{0pt}%
\begin{tabular}{c}#2\end{tabular}}%
}


%%%%%%%%%%%%%%%%%%%%%%%%%%%%%%%%%%%%%%%%%%%%%%%%%%%%%%%%%%%%%
% define the bib style

\usepackage[authoryear]{natbib}
%\usepackage[]{natbib}

%%%%%%%%%%%%%%%%%%%%%%%%%%%%%%%%%%%%%%%%%%%%%%%%%%%%%%%%%%%%%

% adjust maximum number of columns that can be used in a matrix environment

% add matrix cols
%\setcounter{MaxMatrixCols}{30}

%%%%%%%%%%%%%%%%%%%%%%%%%%%%%%%%%%%%%%%%%%%%%%%%%%%%%%%%%%%%%
% all about spaces

\raggedbottom % Prevents spreading out vertical space
\parskip=6pt % Fixed space between paragraphs
\baselineskip=12pt % Fixed baseline skip


%%%%%%%%%%%%%%%%%%%%%%%%%%%%%%%%%%%%%%%%%%%%%%%%%%%%%%%%%%%%%
% config the dedication section

\newcommand\dedication[1]{%
	\thispagestyle{empty}%
	\vspace*{\fill}%
	\begin{flushright}%
		#1%
	\end{flushright}%
%	\vspace*{\fill}%
%	\cleardoublepage%
}
\newcommand\epigraph[3][2\baselineskip]{%
	\begin{flushright}%
		{\rmfamily\itshape #2}%
		\vskip 0.5\baselineskip%
		#3%
	\end{flushright}%
	\vspace{#1}%
}



%%%%%%%%%%%%%%%%%%%%%%%%%%%%%%%%%%%%%%%%%%%%%%%%%%%%%%%%%%%%%%%%

% include other definitions that you have defined yourself
%%%%%%%%%%%%%%%%%%%%%%%%%%%%%%%%%%%%%%%%%%%%%%%%%%%%%%%%%%%%%

% include here your own macros

\usepackage{multirow}
\usepackage{booktabs}

% Pages with this attribute will be displayed in landscape orientation by conforming PDF viewers
\usepackage{pdflscape}
\usepackage{hhline}
\usepackage{enumitem}
%\usepackage{subcaption}
\usepackage{float}
\usepackage{threeparttable}
\usepackage{tikz} % FvW: to make very nice looking graphs
\usepackage{CJKutf8} % for insert Chinese
\usepackage[utf8]{inputenc}
\usepackage{wrapfig}

%%%%%%%%%%%%%%%%%%%%%%%%%%%%%%%%%%%%%%%%%%%%%%%%%%%%%%%%%%%%%%%

% include following package to be able to have line numbers in from of
% all text of thesis. this can be useful to include in the draft
% version of your thesis, such that the committee members can easily
% point out where they have remarks/suggestions/corrections.

\usepackage[left,modulo]{lineno}


\begin{document}

%%%%%%%%%%%%%%%%%%%%%%%%%%%%%%%%%%%%%%%%%%%%%%%%%%%%%%%%%%%%%%%%

\frontmatter

%% define relative path to images for this .tex file
\graphicspath{{./ch-x-others/images/}}


%%%%%%%%%%%%%%%%%%%%%%%%%%%%%%%%%%%%%%%%%%%%%%%%%%%%%%%%%%%%%%%%%%%%
% define below the title of your thesis, the subtitle, and your own name

\newcommand{\thetitle}{Using Social Media to Characterise Crowds in City Events for Crowd Management} 
\newcommand{\thesubtitle}{} 
\newcommand{\theauthor}{Vincent X.\ Gong}


%%%%%%%%%%%%%%%%%%%%%%%%%%%%%%%%%%%%%%%%%%%%%%%%%%%%%%%%%%%%%%%%%%%%
%% page 1: The initial page with only the title and the author name
%%%%%%%%%%%%%%%%%%%%%%%%%%%%%%%%%%%%%%%%%%%%%%%%%%%%%%%%%%%%%%%%%%%%
\pagestyle{fancy} %% start page headers and footers
\thispagestyle{empty}
\vspace*{30mm}

\begin{center}
    \textbf{\huge \thetitle}\\[\baselineskip]\textbf{\LARGE \thesubtitle}
\end{center}


\vspace{40mm}

\begin{center}
    \Large\theauthor
\end{center}

%%%%%%%%%%%%%%%%%%%%%%%%%%%%%%%%%%%%%%%%%%%%%%%%%%%%%%%%%%%%%%%%%%%%
%% page3: the back
%%%%%%%%%%%%%%%%%%%%%%%%%%%%%%%%%%%%%%%%%%%%%%%%%%%%%%%%%%%%%%%%%%%%
\newpage

\thispagestyle{empty}

\vspace*{\fill}
% possibly include your logos and make acknowledgements to funding

\noindent \includegraphics[width=7em]{erc.png}\hspace*{0.75cm}
\includegraphics[width=7em]{allegro.png}\hspace*{0.75cm}
\includegraphics[width=12.5em]{ams.png}\hspace*{0.75cm}\\
%\noindent \includegraphics[width=21em]{logo_allegro}\\

\noindent The research leading to this dissertation has received funding from the European Research Council under the European Union Horizon 2020 Framework Programme for Research and Innovation. It is established by the Scientific Council of the ERC Grant Agreement no. 669792 (Allegro).
%\vspace*{\fill}

% in case there is an illustrator info on the illustrators
%%\vspace*{\fill}
%%\noindent \begin{minipage}[b]{1.0\columnwidth}%
%%\noindent Cover illustration: xxx\\
%%Cover design: xxx \end{minipage}%

%%% if this page should be empty, use the following command to
%%% circumvent de Acrobat error message about wrong contents object
%%% type:
%\vspace*{\fill}
%{\tiny .}


%%%%%%%%%%%%%%%%%%%%%%%%%%%%%%%%%%%%%%%%%%%%%%%%%%%%%%%%%%%%%%%%%%%%%
%% page 3: the page as required by the "pedel"
%%%%%%%%%%%%%%%%%%%%%%%%%%%%%%%%%%%%%%%%%%%%%%%%%%%%%%%%%%%%%%%%%%%%%
\newpage
\thispagestyle{empty}
\vspace*{30mm}

\begin{center}
     \textbf{\huge \thetitle}\\[\baselineskip]\textbf{\LARGE \thesubtitle}
\end{center}

\vspace{20mm}

\begin{center}
    \textbf{Dissertation}
\end{center}
%\vspace{7mm}
\smallskip{}

\selectlanguage{american}

%\noindent{\Huge \textbf{Publications}}
\begin{small}

\begin{center}for the purpose of obtaining the degree of doctor\\
\vspace{1mm}
at Delft University of Technology\\
\vspace{1mm}
by the authority of the Rector Magnificus, Prof. dr. ir. T.H.J.J. van der Hagen\\
\vspace{1mm}
chair of the Board for Doctorates\\
\vspace{1mm}
to be defended publicly on\\
\vspace{1mm}
Monday 21st, September 2020 at 15:00 o'clock
\vspace{1mm}
\vspace{10mm} %5mm vertical space

by

\vspace{10mm}
\textbf{\theauthor}\\
\vspace{1mm}
Master of Science in Computer Science\\
\vspace{1mm}
Delft University of Technology\\
\vspace{1mm}
born in Jiajiang, Basuria
\vspace{1mm}
\end{center}

\end{small}

% ----------------------------------------------
% Start, below is the title page for Graduate School when submitting the Form B.

%\newpage
%\thispagestyle{empty}
%\vspace*{30mm}
%
%\begin{center}
%     \textbf{\huge \thetitle}\\[\baselineskip]\textbf{\LARGE \thesubtitle}
%\end{center}
%%\vspace{30mm}
%%\noindent \begin{center}\begin{tabbing}
%%\hspace*{70mm}\=\kill
%%\noindent Independent members:\\
%%Prof.dr.\ E.\ Simperl \> King's College London\\
%%Dr.\ M.\ Gonzalez \> University of California, Berkeley\\
%%Prof.dr.lr\ M.\ Van\ Steen \> University of Twente\\
%%Prof.dr.ir\ G.J.P.M.\ Houben \> Delft University of Technology\\
%%\end{tabbing}\end{center}
%
%\begin{center}
%    \textbf{Independent members}
%\end{center}\vspace{3mm}
%
%\begin{center}
%Prof.dr. E. Simperl, King's College London\\
%Dr. M. Gonzalez, University of California, Berkeley\\
%Prof.dr.lr M. Van Steen, University of Twente\\
%Prof.dr.ir G.J.P.M. Houben, Delft University of Technology\\
%\end{center}\vspace{3mm}
%
%\begin{center}
%    \textbf{Dissertation}
%\end{center}\vspace{3mm}
%
%%\smallskip{}
%
%\selectlanguage{american}
%
%%\noindent{\Huge \textbf{Publications}}
%\begin{small}
%
%\begin{center}for the purpose of obtaining the degree of doctor\\
%\vspace{1mm}
%at Delft University of Technology\\
%\vspace{1mm}
%by the authority of the Rector Magnificus, Prof.dr.ir. T.H.J.J. van den Hagen\\
%\vspace{1mm}
%chair of the Board for Doctorates\\
%\vspace{1mm}
%to be defended publicly on\\
%\vspace{1mm}
%Monday xx, Month 2020 at 00:00 o'clock
%\vspace{1mm}
%\vspace{4mm} %5mm vertical space
%
%by
%
%\vspace{4mm}
%\textbf{Vincent X. Gong}\\
%\vspace{1mm}
%Master of Science in Computer Science\\
%\vspace{1mm}
%Delft University of Technology\\
%\vspace{1mm}
%born in Jiajiang, Sichuan
%\vspace{1mm}
%\end{center}
%\end{small}

% end, the title page for Graduate School when submitting the Form B.
% ----------------------------------------------

%%%%%%%%%%%%%%%%%%%%%%%%%%%%%%%%%%%%%%%%%%%%%%%%%%%%%%%%%%%%%%%%%%%%%%%%%%%%%%%%%%%%%%%%%
%% page 4: an overview of your review committee
%%%%%%%%%%%%%%%%%%%%%%%%%%%%%%%%%%%%%%%%%%%%%%%%%%%%%%%%%%%%%%%%%%%%%%%%%%%%%%%%%%%%%%%%%
\newpage %
\thispagestyle{empty}

\begin{small}

\noindent This dissertation has been approved by the promotors.\\
\bigskip{}

\noindent Composition of the doctoral committee:
\smallskip{}

\noindent \begin{flushleft}\begin{tabbing}
\hspace*{70mm}\=\kill

Rector Magnificus \> Chairperson\\
Prof. dr. ir.\ S. P. Hoogendoorn \> Delft University of Technology, promotor\\
Dr. ir.\ W. Daamen \> Technische Universiteit Delft, promotor\\
Prof. dr. ir.\ A. Bozzon \> Technische Universiteit Delft, copromotor\\
\bigskip{}\\

\noindent Independent members:\\
\\

Prof. dr. ir.\ G. J. P. M.\ Houben \> Delft University of Technology\\
Prof. dr.\ M.\ Brambilla \> Politecnico di Milano, Italy\\
Dr.\ M.\ Gonzalez \> University of California, Berkeley, U.S.\\
Prof. dr. ir.\ M.\ Van\ Steen \> University of Twente\\
\end{tabbing}\end{flushleft}

\vspace*{\fill}
\selectlanguage{american}

% if you are publishing in the trail series, enable the following with the correct series number
%\noindent TRAIL Thesis Series T2007/14, The Netherlands TRAIL Research School

\bigskip{}
\noindent\textbf{TRAIL Thesis Series no. T2020/xx, The Netherlands Research School TRAIL}\\
\smallskip{}

\noindent TRAIL\\
P.O. BOX 5017\\
2600 GA Delft\\
The Netherlands\\
E-mail: info@rsTRAIL.nl\\
\bigskip{}

\noindent ISBN: your-isbn-number\\
\smallskip{}

%\noindent Copyright \copyright~2020 by \theauthor\\
\noindent Copyright \copyright~2020 by \theauthor\\
\smallskip{}

\noindent All rights reserved. No part of the material protected by this copyright notice may be reproduced or utilized in any form or by any means, electronic or mechanical, including photocopying, recording or by any information storage and retrieval system, without written permission of the author.
\\
% \smallskip{}

\noindent Printed in the Netherlands

\end{small}

}
% define relative path to images for this .tex file
\graphicspath{{./0-misc/images/}}


%%%%%%%%%%%%%%%%%%%%%%%%%%%%%%%%%%%%%%%%%%%%%%%%%%%%%%%%%%%%%%%%%%%%
% define below the title of your thesis, the subtitle, and your own name

\newcommand{\thetitle}{Using Social Media to Characterise Crowds in City Events for Crowd Management} 
\newcommand{\thesubtitle}{} 
\newcommand{\theauthor}{Your\ Name}


%%%%%%%%%%%%%%%%%%%%%%%%%%%%%%%%%%%%%%%%%%%%%%%%%%%%%%%%%%%%%%%%%%%%
%% First the initial page with only the title and the author name
%%%%%%%%%%%%%%%%%%%%%%%%%%%%%%%%%%%%%%%%%%%%%%%%%%%%%%%%%%%%%%%%%%%%
\pagestyle{fancy} %% start page headers and footers
\thispagestyle{empty}
\vspace*{30mm}

\begin{center}
    \textbf{\huge \thetitle}\\[\baselineskip]\textbf{\LARGE \thesubtitle}
\end{center}


\vspace{40mm}

\begin{center}
    \Large\theauthor
\end{center}

%%%%%%%%%%%%%%%%%%%%%%%%%%%%%%%%%%%%%%%%%%%%%%%%%%%%%%%%%%%%%%%%%%%%
%% the back (page 2)
%%%%%%%%%%%%%%%%%%%%%%%%%%%%%%%%%%%%%%%%%%%%%%%%%%%%%%%%%%%%%%%%%%%%
\newpage

\thispagestyle{empty}

\vspace*{\fill}
% possibly include your logos and make acknowledgements to funding

\noindent \includegraphics[width=7em]{erc.png}\hspace*{0.75cm}
\includegraphics[width=7em]{allegro.png}\hspace*{0.75cm}
\includegraphics[width=12.5em]{ams.png}\hspace*{0.75cm}\\
%\noindent \includegraphics[width=21em]{logo_allegro}\\

\noindent The research leading to this dissertation has received funding from the European Research Council under the European Union Horizon 2020 Framework Programme for Research and Innovation. It is established by the Scientific Council of the ERC Grant Agreement no. 669792 (Allegro).
%\vspace*{\fill}

% in case there is an illustrator info on the illustrators
%%\vspace*{\fill}
%%\noindent \begin{minipage}[b]{1.0\columnwidth}%
%%\noindent Cover illustration: xxx\\
%%Cover design: xxx \end{minipage}%

%%% if this page should be empty, use the following command to
%%% circumvent de Acrobat error message about wrong contents object
%%% type:
%\vspace*{\fill}
%{\tiny .}


%%%%%%%%%%%%%%%%%%%%%%%%%%%%%%%%%%%%%%%%%%%%%%%%%%%%%%%%%%%%%%%%%%%%%
%% on page 3 the page as required by the "pedel"
%%%%%%%%%%%%%%%%%%%%%%%%%%%%%%%%%%%%%%%%%%%%%%%%%%%%%%%%%%%%%%%%%%%%%
\newpage
\thispagestyle{empty}
\vspace*{30mm}

\begin{center}
     \textbf{\huge \thetitle}\\[\baselineskip]\textbf{\LARGE \thesubtitle}
\end{center}

\vspace{20mm}

\begin{center}
    \textbf{Dissertation}
\end{center}
%\vspace{7mm}
\smallskip{}

\selectlanguage{american}

%\noindent{\Huge \textbf{Publications}}
\begin{small}

\begin{center}for the purpose of obtaining the degree of doctor\\
\vspace{1mm}
at Delft University of Technology\\
\vspace{1mm}
by the authority of the Rector Magnificus, Prof. dr. ir. T.H.J.J. van der Hagen\\
\vspace{1mm}
chair of the Board for Doctorates\\
\vspace{1mm}
to be defended publicly on\\
\vspace{1mm}
Monday 21st, September 2020 at 15:00 o'clock
\vspace{1mm}
\vspace{10mm} %5mm vertical space

by

\vspace{10mm}
\textbf{\theauthor}\\
\vspace{1mm}
Master of Science in Computer Science\\
\vspace{1mm}
Delft University of Technology\\
\vspace{1mm}
born in Jiajiang, Sichuan
\vspace{1mm}
\end{center}

\end{small}

% ----------------------------------------------
% Start, below is the title page for Graduate School when submitting the Form B.

%\newpage
%\thispagestyle{empty}
%\vspace*{30mm}
%
%\begin{center}
%     \textbf{\huge \thetitle}\\[\baselineskip]\textbf{\LARGE \thesubtitle}
%\end{center}
%%\vspace{30mm}
%%\noindent \begin{center}\begin{tabbing}
%%\hspace*{70mm}\=\kill
%%\noindent Independent members:\\
%%Prof.dr.\ E.\ Simperl \> King's College London\\
%%Dr.\ M.\ Gonzalez \> University of California, Berkeley\\
%%Prof.dr.lr\ M.\ Van\ Steen \> University of Twente\\
%%Prof.dr.ir\ G.J.P.M.\ Houben \> Delft University of Technology\\
%%\end{tabbing}\end{center}
%
%\begin{center}
%    \textbf{Independent members}
%\end{center}\vspace{3mm}
%
%\begin{center}
%Prof.dr. E. Simperl, King's College London\\
%Dr. M. Gonzalez, University of California, Berkeley\\
%Prof.dr.lr M. Van Steen, University of Twente\\
%Prof.dr.ir G.J.P.M. Houben, Delft University of Technology\\
%\end{center}\vspace{3mm}
%
%\begin{center}
%    \textbf{Dissertation}
%\end{center}\vspace{3mm}
%
%%\smallskip{}
%
%\selectlanguage{american}
%
%%\noindent{\Huge \textbf{Publications}}
%\begin{small}
%
%\begin{center}for the purpose of obtaining the degree of doctor\\
%\vspace{1mm}
%at Delft University of Technology\\
%\vspace{1mm}
%by the authority of the Rector Magnificus, Prof.dr.ir. T.H.J.J. van den Hagen\\
%\vspace{1mm}
%chair of the Board for Doctorates\\
%\vspace{1mm}
%to be defended publicly on\\
%\vspace{1mm}
%Monday xx, Month 2020 at 00:00 o'clock
%\vspace{1mm}
%\vspace{4mm} %5mm vertical space
%
%by
%
%\vspace{4mm}
%\textbf{Vincent X. Gong}\\
%\vspace{1mm}
%Master of Science in Computer Science\\
%\vspace{1mm}
%Delft University of Technology\\
%\vspace{1mm}
%born in Jiajiang, Sichuan
%\vspace{1mm}
%\end{center}
%\end{small}

% end, the title page for Graduate School when submitting the Form B.
% ----------------------------------------------

%%%%%%%%%%%%%%%%%%%%%%%%%%%%%%%%%%%%%%%%%%%%%%%%%%%%%%%%%%%%%%%%%%%%%%%%%%%%%%%%%%%%%%%%%
%% and on page 4 an overview of your review committee
%%%%%%%%%%%%%%%%%%%%%%%%%%%%%%%%%%%%%%%%%%%%%%%%%%%%%%%%%%%%%%%%%%%%%%%%%%%%%%%%%%%%%%%%%
\newpage %
\thispagestyle{empty}

\begin{small}

\noindent This dissertation has been approved by the promotors.\\
\bigskip{}

\noindent Composition of the doctoral committee:
\smallskip{}

\noindent \begin{flushleft}\begin{tabbing}
\hspace*{70mm}\=\kill

Rector Magnificus \> Chairperson\\
% rem blow text until GS asks
Prof. dr. ir.\ S. P. Hoogendoorn \> Delft University of Technology, promotor\\
Dr. ir.\ W. Daamen \> Technische Universiteit Delft, promotor\\
Prof. dr. ir.\ A. Bozzon \> Technische Universiteit Delft, copromotor\\
% rem above text until GS asks
\bigskip{}\\

\noindent Independent members:\\
\\

% rem blow text until GS asks
Prof. dr. ir.\ G. J. P. M.\ Houben \> Delft University of Technology\\
Prof. dr.\ M.\ Brambilla \> Politecnico di Milano, Italy\\
Dr.\ M.\ Gonzalez \> University of California, Berkeley, U.S.\\
Prof. dr. ir.\ M.\ Van\ Steen \> University of Twente\\
% rem above text until GS asks
\end{tabbing}\end{flushleft}

\vspace*{\fill}
\selectlanguage{american}

% if you are publishing in the trail series, uncomment the following and include your number
%\noindent TRAIL Thesis Series T2007/14, The Netherlands TRAIL Research School

\bigskip{}
\noindent\textbf{TRAIL Thesis Series no. T2020/xx, the Netherlands Research School TRAIL}\\
\smallskip{}

\noindent TRAIL\\
P.O. BOX 5017\\
2600 GA Delft\\
The Netherlands\\
E-mail: info@rsTRAIL.nl\\
\bigskip{}

\noindent ISBN: your-isbn-number\\
\smallskip{}

%\noindent Copyright \copyright~2020 by \theauthor\\
\noindent Copyright \copyright~2020 by \theauthor\\
\smallskip{}

\noindent All rights reserved. No part of the material protected by this copyright notice may be reproduced or utilized in any form or by any means, electronic or mechanical, including photocopying, recording or by any information storage and retrieval system, without written permission of the author.
\\
% \smallskip{}

\noindent Printed in the Netherlands

\end{small}

%%%%%%%%%%%%%%%%%%%%%%%%%%%%%%%%%%%%%%%%%%%%%%%%%%%%%%%%%%%%%%%%%%%%
%% the back (page 5)
%%%%%%%%%%%%%%%%%%%%%%%%%%%%%%%%%%%%%%%%%%%%%%%%%%%%%%%%%%%%%%%%%%%%
%\newpage
%\thispagestyle{empty}
%
%\vspace*{\fill}
%\rightline{I will say something.}
%\vspace*{\fill}

% in case there is an illustrator info on the illustrators
%%\vspace*{\fill}
%%\noindent \begin{minipage}[b]{1.0\columnwidth}%
%%\noindent Cover illustration: xxx\\
%%Cover design: xxx \end{minipage}%

%%% if this page should be empty, use the following command to
%%% circumvent de Acrobat error message about wrong contents object
%%% type:
%\vspace*{\fill}
%{\tiny .}



%%% Local Variables: 
%%% mode: latex
%%% TeX-master: "../2007-thesis"
%%% End: 

\clearemptydoublepage

%%%%%%%%%%%%%%%%%%%%%%%%%%%%%%%%%%%%%%%%%%%%%%%%%%%%%%%%%%%%%%%%
\mainmatter

% here comes the preface

% to make sure that from here on american hyphenation is used, include
% the following

\selectlanguage{american}

\pagestyle{fancy}
\pagenumbering{roman}

% the 5 title pages defined in 0-misc/title.text have no numbers,
% so we manually set the next page, which should have a page number, to 6

\setcounter{page}{5}
%%%%%%%%%%%%%%%%%%%%%%%%%%%%%%%%%%%%%%%%%%%%%%%%%%%%%%%%%%%%%%%%

\clearemptydoublepage
%\dedication{\epigraph{I have something to say.}{}}
%\vspace*{-0.4cm}
%\vspace*{4cm}
\dedication{\epigraph{I have something to say.}{X. Gong}}



\clearemptydoublepage

%%%%%%%%%%%%%%%%%%%%%%%%%%%%%%%%%%%%%%%%%%%%%%%%%%%%%%%%%%%%%%%%


\chapter*{Preface}
\addtocontents{toc}{\vspace*{-0.9cm}}
\addcontentsline{toc}{chapter}{Preface}

\graphicspath{{./ch-x-others/images/}}

\vspace*{-0.4cm}


Preface
\clearemptydoublepage

%%%%%%%%%%%%%%%%%%%%%%%%%%%%%%%%%%%%%%%%%%%%%%%%%%%%%%%%%%%%%%%%

% here comes the table of contents

\tableofcontents
\clearemptydoublepage

%%%%%%%%%%%%%%%%%%%%%%%%%%%%%%%%%%%%%%%%%%%%%%%%%%%%%%%%%%%%%%%%

%% reset numbering and set numbering to arabic
\pagenumbering{arabic}
\setcounter{page}{1}

%%%%%%%%%%%%%%%%%%%%%%%%%%%%%%%%%%%%%%%%%%%%%%%%%%%%%%%%%%%%%%%%

% here come the main chapters

% for reviewing purposes it may be useful to include line numbers.
% (this requires \usepackage{lineno} above to be uncommented)
\begin{linenumbers}[1]

% chapter 1
\include{1-intro/intro}
\clearemptydoublepage 

% chapter 2
\chapter{Introduction}\label{ch1_title}
% the short chapter mark for page header 
\chaptermark{Introduction}

\graphicspath{{./ch-1-intro/images/}}

City-scale events are getting more popular and attract a large number of people participating in various activities. For instance, on King’s Day, a national holiday in the Netherlands, a huge amount of people pour into the city and gather in the urban area, participating in various activities such as street parties, music festivals and boat parades. Event stakeholders, such as event organisers, police, municipalities, and crowd managers manage the crowd to avoid incidents. Crowd management practice consists of two phases \citep{martella2017current}, i.e. the planning phase and operational phase. In the planning phase, crowd managers require the past event data to infer guidelines, and to perform computer simulations of the crowds in the event.


Parts of this chapter have been published in \cite{NegDeS:05-021}.

\section{Background}\label{ch1_background}
Some background about this research.

\clearemptydoublepage

\include{6-conc/conc}

%\clearemptydoublepage


%%%%%%%%%%%%%%%%%%%%%%%%%%%%%%%%%%%%%%%%%%%%%%%%%%%%%%%%%%%%%%%%

% here come the appendices, if any

%\appendix

%\input{\fname{A1-my_appendix/my_appendix}}
%\include{A1-my_appendix/my_appendix}
%\clearemptydoublepage


%%%%%%%%%%%%%%%%%%%%%%%%%%%%%%%%%%%%%%%%%%%%%%%%%%%%%%%%%%%%%%%%

% here comes the bibliography

%% next line is necessary to make link in table of contents appear correctly
\phantomsection
\addcontentsline{toc}{chapter}{Bibliography}

% add the bibliography file, by default it is the references.bib
\bibliography{./ch-x-others/references}
}
% next line is necessary to make link in table of contents appear correctly
\phantomsection
\addcontentsline{toc}{chapter}{Bibliography}

% add the bibliography file, by default it is the references.bib
\bibliography{./ch-x-others/references}

\clearemptydoublepage

%%%%%%%%%%%%%%%%%%%%%%%%%%%%%%%%%%%%%%%%%%%%%%%%%%%%%%%%%%%%%%%%

% here comes the glossary (partially constructed by the commands in
% 0-misc/lists.tex, which was included above)

%\chapter*{Glossary}\markboth{Glossary}{Glossary}
\addcontentsline{toc}{chapter}{Glossary}

\graphicspath{{./ch-x-others/images/}}


\section*{Conventions}
The following conventions are used in this thesis for notation
and symbols:
\begin{itemize}
\item A lower case character typeset in boldface, e.g., $\mathbf{x}$,
   represents a column vector.
\item The number of elements in a vector $\mathbf{x}$ is indicated by
   $n_{\mathbf{x}}$.
\item \ldots
\end{itemize}

\section*{List of symbols and notations}
Below follows a list of the most frequently used symbols and notations
in this thesis. Symbols particular to power network applications
are explained only in the relevant chapters.

\vspace*{-0.4cm}\begin{tabbing}
\hspace*{3.5cm}\=\kill\\
$ \mathbf{A}  $ \> system matrices of linear time-invariant models \\
\\[-0.1cm]
$ \mathbf{B}  $ \> input matrices of linear time-invariant models \\
\\[-0.1cm]
$ \mathbf{E}_1 $ \> matrices of mixed-logical dynamic models \\
\end{tabbing}

\vspace*{-0.7cm}
\section*{List of abbreviations}
The following abbreviations are used in this thesis:
\begin{tabbing}
\hspace*{4cm}\=\kill
AVR \> Automatic Voltage Regulator\\
DAE \> Differential-Algebraic Equations\\
FACTS \> Flexible Alternating-Current Transmission System\\
MPC \> Model Predictive Control\\
\ldots \> \ldots
\end{tabbing}
}
\chapter*{Glossary}\markboth{Glossary}{Glossary}
\addcontentsline{toc}{chapter}{Glossary}

\graphicspath{{./ch-x-others/images/}}


\section*{Conventions}
The following conventions are used in this thesis for notation
and symbols:
\begin{itemize}
\item A lower case character typeset in boldface, e.g., $\mathbf{x}$,
   represents a column vector.
\item The number of elements in a vector $\mathbf{x}$ is indicated by
   $n_{\mathbf{x}}$.
\item \ldots
\end{itemize}

\section*{List of symbols and notations}
Below follows a list of the most frequently used symbols and notations
in this thesis. Symbols particular to power network applications
are explained only in the relevant chapters.

\vspace*{-0.4cm}\begin{tabbing}
\hspace*{3.5cm}\=\kill\\
$ \mathbf{A}  $ \> system matrices of linear time-invariant models \\
\\[-0.1cm]
$ \mathbf{B}  $ \> input matrices of linear time-invariant models \\
\\[-0.1cm]
$ \mathbf{E}_1 $ \> matrices of mixed-logical dynamic models \\
\end{tabbing}

\vspace*{-0.7cm}
\section*{List of abbreviations}
The following abbreviations are used in this thesis:
\begin{tabbing}
\hspace*{4cm}\=\kill
AVR \> Automatic Voltage Regulator\\
DAE \> Differential-Algebraic Equations\\
FACTS \> Flexible Alternating-Current Transmission System\\
MPC \> Model Predictive Control\\
\ldots \> \ldots
\end{tabbing}

\clearemptydoublepage
%%%%%%%%%%%%%%%%%%%%%%%%%%%%%%%%%%%%%%%%%%%%%%%%%%%%%%%%%%%%%%%%
% here comes the summary in english

%\chapter*{Summary}\markboth{Summary}{Summary}
\addcontentsline{toc}{chapter}{Summary}

\graphicspath{{./ch-x-others/images/}}

Events are getting more popular and more frequent in cities around the world. In the Netherlands in 2017, the number of festivals grew to almost 1000 \footnote{https://www.eventbranche.nl/nieuws/aantal-festivals-groeit-tot-bijna-1000-per-jaar-aantal-bezoeken-daalt-miniem-16483.html}. These events take place in large areas of the city, they have a common topic, they include sub-events (activities), and they have start and end times and lasts from one day to several days. Examples of events are the national holidays, Soul Live Festival and trade exhibitions. City events can easily attract a large number of people. Event stakeholders, such as the event organizers, police, municipalities and other authorities, and crowd managers are concerned with guaranteeing the safety, comfort and general well being of the attendees. To this end, they enforce predefined crowd management measures that are adaptive to the current state of the event environment and of the participating crowd. This state is measured through information about the factors influencing event planning \citep{li2019crowds} and pedestrian behaviour \citep{still2000crowd,tubbs2007egress,abbott2000event, zomer2015managing} for crowd management, such as crowd size, density, mobility, emotion, visitor profile, and location. Conventionally, this information is derived from data provided by stewards (operating on the ground during the event) and sometimes pre-installed sensing infrastructures, such as counting systems, Bluetooth/ Wi-Fi sensors, and video cameras. While effective, these solutions suffer from several issues: they provide little information about sentiments, gender and age distribution, they are expensive, they cannot provide Spatio-temporal information, and they are complex to install and maintain.  

}
\chapter*{Summary}\markboth{Summary}{Summary}
\addcontentsline{toc}{chapter}{Summary}

\graphicspath{{./ch-x-others/images/}}

Events are getting more popular and more frequent in cities around the world. In the Netherlands in 2017, the number of festivals grew to almost 1000 \footnote{https://www.eventbranche.nl/nieuws/aantal-festivals-groeit-tot-bijna-1000-per-jaar-aantal-bezoeken-daalt-miniem-16483.html}. These events take place in large areas of the city, they have a common topic, they include sub-events (activities), and they have start and end times and lasts from one day to several days. Examples of events are the national holidays, Soul Live Festival and trade exhibitions. City events can easily attract a large number of people. Event stakeholders, such as the event organizers, police, municipalities and other authorities, and crowd managers are concerned with guaranteeing the safety, comfort and general well being of the attendees. To this end, they enforce predefined crowd management measures that are adaptive to the current state of the event environment and of the participating crowd. This state is measured through information about the factors influencing event planning \citep{li2019crowds} and pedestrian behaviour \citep{still2000crowd,tubbs2007egress,abbott2000event, zomer2015managing} for crowd management, such as crowd size, density, mobility, emotion, visitor profile, and location. Conventionally, this information is derived from data provided by stewards (operating on the ground during the event) and sometimes pre-installed sensing infrastructures, such as counting systems, Bluetooth/ Wi-Fi sensors, and video cameras. While effective, these solutions suffer from several issues: they provide little information about sentiments, gender and age distribution, they are expensive, they cannot provide Spatio-temporal information, and they are complex to install and maintain.  


\clearemptydoublepage

%%%%%%%%%%%%%%%%%%%%%%%%%%%%%%%%%%%%%%%%%%%%%%%%%%%%%%%%%%%%%%%%
% here comes a summary in dutch

%\chapter*{Samenvatting}\markboth{Samenvatting}{Samenvatting}
\addcontentsline{toc}{chapter}{Samenvatting (Summary in Dutch)}

\graphicspath{{./ch-x-others/images/}}

% Samenvatting should be in dutch, so use dutch hyphenation
\selectlanguage{dutch}

\section*{Social media gebruiken om menigte te karakteriseren in stadsevenementen voor crowd management}

Een samenvatting in het Nederlands zal hier worden gepresenteerd.


\bigskip
\noindent Vincent X.\ Gong

\selectlanguage{american}
}
\chapter*{Samenvatting}\markboth{Samenvatting}{Samenvatting}
\addcontentsline{toc}{chapter}{Samenvatting (Summary in Dutch)}

\graphicspath{{./ch-x-others/images/}}

% Samenvatting should be in dutch, so use dutch hyphenation
\selectlanguage{dutch}

\section*{Social media gebruiken om menigte te karakteriseren in stadsevenementen voor crowd management}

Een samenvatting in het Nederlands zal hier worden gepresenteerd.


\bigskip
\noindent Vincent X.\ Gong

\selectlanguage{american}

\clearemptydoublepage
%%%%%%%%%%%%%%%%%%%%%%%%%%%%%%%%%%%%%%%%%%%%%%%%%%%%%%%%%%%%%%%%
% here comes the summary in other languages.
% you can remove it if you do not want it.

%\chapter*{Summary}\markboth{Summary}{Summary}
\addcontentsline{toc}{chapter}{Summary}

\graphicspath{{./ch-x-others/images/}}

Events are getting more popular and more frequent in cities around the world. In the Netherlands in 2017, the number of festivals grew to almost 1000 \footnote{https://www.eventbranche.nl/nieuws/aantal-festivals-groeit-tot-bijna-1000-per-jaar-aantal-bezoeken-daalt-miniem-16483.html}. These events take place in large areas of the city, they have a common topic, they include sub-events (activities), and they have start and end times and lasts from one day to several days. Examples of events are the national holidays, Soul Live Festival and trade exhibitions. City events can easily attract a large number of people. Event stakeholders, such as the event organizers, police, municipalities and other authorities, and crowd managers are concerned with guaranteeing the safety, comfort and general well being of the attendees. To this end, they enforce predefined crowd management measures that are adaptive to the current state of the event environment and of the participating crowd. This state is measured through information about the factors influencing event planning \citep{li2019crowds} and pedestrian behaviour \citep{still2000crowd,tubbs2007egress,abbott2000event, zomer2015managing} for crowd management, such as crowd size, density, mobility, emotion, visitor profile, and location. Conventionally, this information is derived from data provided by stewards (operating on the ground during the event) and sometimes pre-installed sensing infrastructures, such as counting systems, Bluetooth/ Wi-Fi sensors, and video cameras. While effective, these solutions suffer from several issues: they provide little information about sentiments, gender and age distribution, they are expensive, they cannot provide Spatio-temporal information, and they are complex to install and maintain.  

}
% \chapter*{Summary in Chinese}\markboth{Summarycn}{Summarycn}
\addcontentsline{toc}{chapter}{Summary in Chinese}

\graphicspath{{./ch-x-others/images/}}


\begin{CJK*}{UTF8}{gbsn}
A summary in Chinese will be presented here.

摘要。
\end{CJK*}

% \clearemptydoublepage

%%%%%%%%%%%%%%%%%%%%%%%%%%%%%%%%%%%%%%%%%%%%%%%%%%%%%%%%%%%%%%%%

% and here comes the curriculum vitae of the author

%\chapter*{About the author}\markboth{About the author}{About the author}

\addcontentsline{toc}{chapter}{About the author}

\begin{wrapfigure}{O}{0.34\textwidth}
%\centering
\includegraphics[width=0.34\textwidth, trim=0 0 0 0.8cm]{./0-misc/images/vincent.jpg}\par
\caption{\label{fig:frog2}Vincent X. Gong}
\end{wrapfigure}

Sed ut perspiciatis unde omnis iste natus error sit voluptatem accusantium doloremque laudantium, totam rem aperiam, eaque ipsa quae ab illo inventore veritatis et quasi architecto beatae vitae dicta sunt explicabo. Nemo enim ipsam voluptatem quia voluptas sit aspernatur aut odit aut fugit, sed quia consequuntur magni dolores eos qui ratione voluptatem sequi nesciunt. Neque porro quisquam est, qui dolorem ipsum quia dolor sit amet, consectetur, adipisci velit, sed quia non numquam eius modi tempora incidunt ut labore et dolore magnam aliquam quaerat voluptatem. Ut enim ad minima veniam, quis nostrum exercitationem ullam corporis suscipit laboriosam, nisi ut aliquid ex ea commodi consequatur? Quis autem vel eum iure reprehenderit qui in ea voluptate velit esse quam nihil molestiae consequatur, vel illum qui dolorem eum fugiat quo voluptas nulla pariatur.\\

At vero eos et accusamus et iusto odio dignissimos ducimus qui blanditiis praesentium voluptatum deleniti atque corrupti quos dolores et quas molestias excepturi sint occaecati cupiditate non provident, similique sunt in culpa qui officia deserunt mollitia animi, id est laborum et dolorum fuga. Et harum quidem rerum facilis est et expedita distinctio.\\


Nam libero tempore, cum soluta nobis est eligendi optio cumque nihil impedit quo minus id quod maxime placeat facere possimus, omnis voluptas assumenda est, omnis dolor repellendus. Temporibus autem quibusdam et aut officiis debitis aut rerum necessitatibus saepe eveniet ut et voluptates repudiandae sint et molestiae non recusandae.


\vspace*{\fill}
\pagebreak

\vspace*{0.60cm}
\noindent{\Huge \textbf{Publications}}

\section*{Journal papers}

\begin{enumerate}
    \item \textbf{Gong, Vincent X.}, et al. "Using social media for attendees density estimation in city-scale events." IEEE Access 6 (2018): 36325-36340.
    \item \textbf{Gong, Vincent X.}, Jie Yang, Winnie Daamen, Alessandro Bozzon, Serge P. Hoogendoorn, and Geert-Jan Houben. "Using social media for attendees density estimation in city-scale events." IEEE Access 6 (2018): 36325-36340.
    
    % \item \textbf{Gong, Vincent X.}, Jie Yang, Winnie Daamen, Alessandro Bozzon, Serge P. Hoogendoorn, and Geert-Jan Houben. "Using social media for attendees density estimation in city-scale events." IEEE Access 6 (2018): 36325-36340.
    
    % \item \textbf{Gong, Vincent X.}, Jie Yang, Winnie Daamen, Alessandro Bozzon, Serge P. Hoogendoorn, and Geert-Jan Houben. "Using social media for attendees density estimation in city-scale events." IEEE Access 6 (2018): 36325-36340.
    
    % \item \textbf{Gong, Vincent X.}, Jie Yang, Winnie Daamen, Alessandro Bozzon, Serge P. Hoogendoorn, and Geert-Jan Houben. "Using social media for attendees density estimation in city-scale events." IEEE Access 6 (2018): 36325-36340.
\end{enumerate}

% to get page numbers a multiple of 8 it may be necessary to add some empty pages

%\pagestyle{plain}\fancyhf{}\markboth{}{}
%\newpage{\mbox{}}

%%% Local Variables: 
%%% mode: latex
%%% TeX-master: "../2007-thesis"
%%% End: 
}
\chapter*{About the author}\markboth{About the author}{About the author}

\addcontentsline{toc}{chapter}{About the author}

\begin{wrapfigure}{O}{0.34\textwidth}
%\centering
\includegraphics[width=0.34\textwidth, trim=0 0 0 0.8cm]{./0-misc/images/vincent.jpg}\par
\caption{\label{fig:frog2}Vincent X. Gong}
\end{wrapfigure}

Sed ut perspiciatis unde omnis iste natus error sit voluptatem accusantium doloremque laudantium, totam rem aperiam, eaque ipsa quae ab illo inventore veritatis et quasi architecto beatae vitae dicta sunt explicabo. Nemo enim ipsam voluptatem quia voluptas sit aspernatur aut odit aut fugit, sed quia consequuntur magni dolores eos qui ratione voluptatem sequi nesciunt. Neque porro quisquam est, qui dolorem ipsum quia dolor sit amet, consectetur, adipisci velit, sed quia non numquam eius modi tempora incidunt ut labore et dolore magnam aliquam quaerat voluptatem. Ut enim ad minima veniam, quis nostrum exercitationem ullam corporis suscipit laboriosam, nisi ut aliquid ex ea commodi consequatur? Quis autem vel eum iure reprehenderit qui in ea voluptate velit esse quam nihil molestiae consequatur, vel illum qui dolorem eum fugiat quo voluptas nulla pariatur.\\

At vero eos et accusamus et iusto odio dignissimos ducimus qui blanditiis praesentium voluptatum deleniti atque corrupti quos dolores et quas molestias excepturi sint occaecati cupiditate non provident, similique sunt in culpa qui officia deserunt mollitia animi, id est laborum et dolorum fuga. Et harum quidem rerum facilis est et expedita distinctio.\\


Nam libero tempore, cum soluta nobis est eligendi optio cumque nihil impedit quo minus id quod maxime placeat facere possimus, omnis voluptas assumenda est, omnis dolor repellendus. Temporibus autem quibusdam et aut officiis debitis aut rerum necessitatibus saepe eveniet ut et voluptates repudiandae sint et molestiae non recusandae.


\vspace*{\fill}
\pagebreak

\vspace*{0.60cm}
\noindent{\Huge \textbf{Publications}}

\section*{Journal papers}

\begin{enumerate}
    \item \textbf{Gong, Vincent X.}, et al. "Using social media for attendees density estimation in city-scale events." IEEE Access 6 (2018): 36325-36340.
    \item \textbf{Gong, Vincent X.}, Jie Yang, Winnie Daamen, Alessandro Bozzon, Serge P. Hoogendoorn, and Geert-Jan Houben. "Using social media for attendees density estimation in city-scale events." IEEE Access 6 (2018): 36325-36340.
    
    % \item \textbf{Gong, Vincent X.}, Jie Yang, Winnie Daamen, Alessandro Bozzon, Serge P. Hoogendoorn, and Geert-Jan Houben. "Using social media for attendees density estimation in city-scale events." IEEE Access 6 (2018): 36325-36340.
    
    % \item \textbf{Gong, Vincent X.}, Jie Yang, Winnie Daamen, Alessandro Bozzon, Serge P. Hoogendoorn, and Geert-Jan Houben. "Using social media for attendees density estimation in city-scale events." IEEE Access 6 (2018): 36325-36340.
    
    % \item \textbf{Gong, Vincent X.}, Jie Yang, Winnie Daamen, Alessandro Bozzon, Serge P. Hoogendoorn, and Geert-Jan Houben. "Using social media for attendees density estimation in city-scale events." IEEE Access 6 (2018): 36325-36340.
\end{enumerate}

% to get page numbers a multiple of 8 it may be necessary to add some empty pages

%\pagestyle{plain}\fancyhf{}\markboth{}{}
%\newpage{\mbox{}}

%%% Local Variables: 
%%% mode: latex
%%% TeX-master: "../2007-thesis"
%%% End: 

%\clearemptydoublepage

%%%%%%%%%%%%%%%%%%%%%%%%%%%%%%%%%%%%%%%%%%%%%%%%%%%%%%%%%%%%%%%%
% here comes a list with trail series phd theses. (no need to include
% these if you are not publishing in the trail series.)

% \chapter*{TRAIL Thesis Series}
\addcontentsline{toc}{chapter}{TRAIL Thesis Series publications}
\markboth{TRAIL Thesis Series}{} \markright{TRAIL Thesis Series}

\noindent The following list contains the most recent dissertations in the TRAIL Thesis Series, For a complete overview of more than 250 titles see the TRAIL website: www.rsTRAIL.nl.

\vspace*{0.40cm}{}
\noindent The TRAIL Thesis Series is a series of the Netherlands TRAIL Research School on transport, infrastructure and logistics.

\vspace*{0.40cm}{}
\noindent Nat, C.G.J.M., van der, \textit{A Knowledge-based Concept
Exploration Model for Submarine Design}, T99/1, March 1999, TRAIL
Thesis Series, Delft University Press, The Netherlands

\vspace*{0.25cm}{} 
\noindent Nuttall, A.J.G., \textit{Design Aspects of Multiple Driven Belt Conveyors}, T2007/12, November 2007, TRAIL Thesis Series, The
Netherlands

\vspace*{0.25cm}{} 
\noindent Nederveen, A.A.J., \textit{Ruimtelijke Inpassing van Lijninfrastructuur. Een onderzoek naar de geschiktheid van inspraakreacties voor het ontwerpen van lijninfrastructuur}, T2007/13, December 2007, TRAIL Thesis Series, The
Netherlands

\vspace*{0.25cm}{} 
\noindent Negenborn, R.R., \textit{Multi-Agent Model Predictive Control with
Applications to Power Networks}, T2007/14, December 2007, TRAIL Thesis
Series, The Netherlands


%%% Local Variables: 
%%% mode: latex
%%% TeX-master: "../2007-thesis"
%%% End: 
}
\chapter*{TRAIL Thesis Series}
\addcontentsline{toc}{chapter}{TRAIL Thesis Series publications}
\markboth{TRAIL Thesis Series}{} \markright{TRAIL Thesis Series}

\noindent The following list contains the most recent dissertations in the TRAIL Thesis Series, For a complete overview of more than 250 titles see the TRAIL website: www.rsTRAIL.nl.

\vspace*{0.40cm}{}
\noindent The TRAIL Thesis Series is a series of the Netherlands TRAIL Research School on transport, infrastructure and logistics.

\vspace*{0.40cm}{}
\noindent Nat, C.G.J.M., van der, \textit{A Knowledge-based Concept
Exploration Model for Submarine Design}, T99/1, March 1999, TRAIL
Thesis Series, Delft University Press, The Netherlands

\vspace*{0.25cm}{} 
\noindent Nuttall, A.J.G., \textit{Design Aspects of Multiple Driven Belt Conveyors}, T2007/12, November 2007, TRAIL Thesis Series, The
Netherlands

\vspace*{0.25cm}{} 
\noindent Nederveen, A.A.J., \textit{Ruimtelijke Inpassing van Lijninfrastructuur. Een onderzoek naar de geschiktheid van inspraakreacties voor het ontwerpen van lijninfrastructuur}, T2007/13, December 2007, TRAIL Thesis Series, The
Netherlands

\vspace*{0.25cm}{} 
\noindent Negenborn, R.R., \textit{Multi-Agent Model Predictive Control with
Applications to Power Networks}, T2007/14, December 2007, TRAIL Thesis
Series, The Netherlands


%%% Local Variables: 
%%% mode: latex
%%% TeX-master: "../2007-thesis"
%%% End: 

%\clearemptydoublepage
%%%%%%%%%%%%%%%%%%%%%%%%%%%%%%%%%%%%%%%%%%%%%%%%%%%%%%%%%%%%%%%%

% that's it!
\end{linenumbers}
\end{document}


%%% Local Variables:
%%% mode: latex
%%% TeX-master: t
%%% End:
